\subsection{The genus count of the original embedded isometries}
\label{sec:embedded-isometry-genus-count}

The fixed-label return bound uses pairs of integral data together
with an isometry compatible with their reductions.  We retain the
actual data and their congruence condition.  The decisive estimate
comes from the determinant of the isometry minus the identity on
the orthogonal complement of its fixed space.

Let \(E/F\) be a CM extension unramified at every finite place,
\([E:\mathbf Q]=2f\), and let \(V\) be a positive definite
Hermitian \(E\)-space of rank \(k\geq1\).  Put
\[
 \ell=\log_2(2f),\qquad c=1+\ell.
\]
Fix an integral lattice \(\Lambda\) with
\(\Delta(\Lambda)=D\), and assume that
\(\Lambda^\#/\Lambda\) is generated by at most one element
over \(\mathcal O_E\).  Let \(g\geq2\) and \(j\geq0\) be integers.

An embedded datum of index \(x\) is a pair \(d=(L,\Gamma)\).
Here \(L\subset V\) is an integral lattice of discriminant \(D\),
and \(\Gamma\subset L\times\Lambda\) is the inverse image of the
graph of an \(\mathcal O_E/g\mathcal O_E\)-linear isometry
\[
                         L/gL\longrightarrow\Lambda/g\Lambda
\]
for their induced Hermitian forms.  Put
\[
 F(d)=\{y\in V:(y,y)\in\Gamma\}.
 \tag{EI1}
\]
We require
\([\Lambda:F(d)]=x\) and that \(\Lambda/F(d)\) have at most
\(j\) generators.  Let \(\mathcal D_j(x;\Lambda,g)\) be the
set of these actual embedded data, without dividing by
isometries.

Let \(\mathcal T_j(x,x';\Lambda,g)\) consist of triples
\((d,d',\eta)\), with \(d=(L,\Gamma)\in\mathcal D_j(x;\Lambda,g)\),
\(d'=(L',\Gamma')\in\mathcal D_j(x';\Lambda,g)\), and an
\(E\)-linear unitary isometry \(\eta\) of \(V\) such that
\[
       \eta L'=L,\qquad (\eta\times1)\Gamma'=\Gamma.
 \tag{EI2}
\]
These are the source's embedded data and triples, with the same
lattice indices, fixed-lattice convention and mod-\(g\) graph.
\cite{embedded-isometry-count-source}
All indices below are cardinalities of additive quotients.

\begin{lemma}[The arithmetic consequence of the congruence]
\label{lem:embedded-isometry-fixed-discriminant}
Let \((d,d',\eta)\) be such a triple.  Set
\[
 U=\ker(\eta-1),\qquad s=k-\dim_EU,\qquad
                  A=F(d)\cap U.
\]
Then \(F(d')\cap U=A\), and
\[
                 \Delta(A)\leq(2/g)^{2sf}Dxx'.
 \tag{EI3}
\]
For \(U=0\), the discriminant on the left is \(1\).
\end{lemma}

\begin{proof}
If \(y\in F(d')\cap U\), then \((y,y)\in\Gamma'\), so
(EI2) gives \((\eta y,y)=(y,y)\in\Gamma\).  Applying the
same argument to \(\eta^{-1}\) proves the equality of the
two intersections.

For \(y\in F(d')\) and \(z\in F(d)\), the two pairs
\((\eta y,y)\) and \((z,z)\) belong to \(\Gamma\).
The mod-\(g\) isometry condition therefore gives
\[
                    h((\eta-1)y,z)\in g\mathcal O_E.
 \tag{EI4}
\]
Put \(W=U^\perp\), and let \(P,P'\) be the orthogonal
projections of \(F(d),F(d')\) onto \(W\).  They are full
fractional lattices in \(W\).  Equation (EI4) becomes
\[
                  (\eta-1)P'\subset gP^\#.
 \tag{EI5}
\]
The map \(\eta-1\) is invertible on \(W\).

For a fractional Hermitian lattice \(B\), write
\(\delta(B)\) for its normalized covolume; when \(B\)
is integral, \(\delta(B)=\sqrt{\Delta(B)}\).
More explicitly, use ordinary real Lebesgue measure in
unitary coordinates at one complex embedding above each
real place, and divide the archimedean covolume by
\[
       \left(2^{-f}\sqrt{|\operatorname{disc}E|}\right)^
                                               {\operatorname{rank}B}.
\]
This is the covolume of the standard \(\mathcal O_E\)-lattice
of the same rank.  Duality gives
\(\delta(B^\#)=\delta(B)^{-1}\), and an \(E\)-linear map
\(a\) changes it by \(|N_{E/\mathbf Q}\det a|\).
The volume formula for an orthogonal projection gives
\[
 \delta(F(d))=\delta(A)\delta(P),\qquad
 \delta(F(d'))=\delta(A)\delta(P').
 \tag{EI6}
\]
These identities may also be checked in local lattice bases;
they do not require globally free lattices.

The inclusion (EI5) has positive integral index.  At every
complex embedding the unitary operator \(\eta\) has
\(\|\eta-1\|\leq2\).  Consequently
\[
 g^{2sf}\delta(P)^{-1}
 \leq |N_{E/\mathbf Q}\det((\eta-1)|_W)|\,\delta(P')
 \leq2^{2sf}\delta(P').
\]
Since
\(\delta(F(d))=\sqrt D\,x\) and
\(\delta(F(d'))=\sqrt D\,x'\), equation (EI6) gives
\(\delta(P)\delta(P')=Dxx'/\Delta(A)\).
This proves (EI3), including the zero-dimensional convention.
\end{proof}

\begin{lemma}[The number of embedded data]
\label{lem:embedded-data-count}
Put \(m=3(2j+1)\).  For a fixed lattice \(B\subset\Lambda\)
of index \(x\geq1\),
\[
 \#\{d\in\mathcal D_j(x;\Lambda,g):F(d)=B\}
                               \leq(xD)^m.
 \tag{EI7}
\]
In particular,
\[
             |\mathcal D_j(x;\Lambda,g)|
                    \leq x^{k+c}(xD)^m.
 \tag{EI8}
\]
\end{lemma}

\begin{proof}
If the fiber is nonempty, \(\Lambda/B\) has at most \(j\)
generators.  The chain
\[
              B\subset\Lambda\subset\Lambda^\#\subset B^\#
\]
shows that \(B^\#/B\) has at most \(2j+1\) generators.
Indeed \(B^\#/\Lambda^\#\) is dual to \(\Lambda/B\);
its local cyclic factors, hence its number of generators,
are the same.  Finite primary components combine by the
Chinese remainder theorem.  The middle quotient has one
generator by hypothesis.

For \(d=(L,\Gamma)\) in the fiber, integrality gives
\(B\subset L\subset B^\#\).  Since \(\Delta(L)=D\),
the discriminant-index identity gives \([L:B]=x\).
The group \(B^\#/B\) has order \(Dx^2\).  Every one of
its submodules has at most \(2j+1\) generators, so there
are at most \((Dx^2)^{2j+1}\) possibilities for \(L\).

For a fixed \(L\), choose for \(y\in\Lambda\) an element
\(y'\in L\) with \((y',y)\in\Gamma\).  Pairing this
element with \((z,z)\in\Gamma\), \(z\in B\), shows
\((y'-y)/g\in B^\#\).  The rule
\[
 \alpha_\Gamma:\Lambda/B\longrightarrow B^\#/L,\qquad
        y+B\longmapsto (y'-y)/g+L
 \tag{EI9}
\]
is an \(\mathcal O_E\)-linear homomorphism.  Changing
\(y'\) changes the quotient by an element of \(L\);
changing \(y\) by an element of \(B\) leaves its class
unchanged.  Conversely \(L\) and \(\alpha_\Gamma\)
determine \(\Gamma\): a pair \((y',y)\in L\times\Lambda\)
belongs to it precisely when \((y'-y)/g\in B^\#\) and
its class in \(B^\#/L\) equals \(\alpha_\Gamma(y+B)\).

The target in (EI9) has \(Dx\) elements, and the domain
has \(j\) generators.  Thus there are at most
\((Dx)^j\) possible graphs for each \(L\).  Their product
is at most
\[
 (Dx^2)^{2j+1}(Dx)^j
       =D^{3j+1}x^{5j+2}\leq(Dx)^{6j+3},
\]
which proves (EI7).  The rank-\(k\) sublattice count in
the proof of (MM13) bounds the number of \(B\)'s by
\(x^{k+\ell}\leq x^{k+c}\), giving (EI8).
\end{proof}

\begin{lemma}[Sublattices with a prescribed intersection]
\label{lem:sublattice-prescribed-intersection}
Let \(U\subset V\) have codimension \(s\), and fix
a lattice \(A\subset\Lambda\cap U\).  The number of
sublattices \(B\subset\Lambda\) with
\([\Lambda:B]=x\geq1\) and \(B\cap U=A\) is at most
\[
                              x^{s+\ell}.
 \tag{EI10}
\]
\end{lemma}

\begin{proof}
The assertion is immediate when the set is empty.
Otherwise put \(A_0=\Lambda\cap U\) and
\(a=[A_0:A]\).  The quotient \(\Lambda/A_0\) is a
projective \(\mathcal O_E\)-module of rank \(s\), so
choose a module splitting \(\Lambda=A_0\oplus C\).
The image \(C'\) of \(B\) in \(C\) has index \(y=x/a\).
For fixed \(C'\), the lattice \(B\) is determined by a
homomorphism \(C'\to A_0/A\).  There are exactly \(a^s\)
such homomorphisms: locally \(C'\) is free of rank \(s\),
and the finite primary components combine independently.
There are at most \(y^{s+\ell}\) choices for \(C'\) by
the rank-\(s\) index count.  Their product is at most
\((ay)^{s+\ell}=x^{s+\ell}\).  For \(s=0\), the only
possible \(B\) is \(A\), and the same bound holds.
\end{proof}

\paragraph{The finite isometry bound.}
For two lattices and a subspace \(U\) of codimension \(s\),
the isometries carrying one lattice to the other and fixing
\(U\) pointwise form either the empty set or a coset of a
finite lattice automorphism group.  Restriction to \(U^\perp\)
embeds that group in \(GL_{2sf}(\mathbf Q)\).
Minkowski's finite-group theorem implies that a finite
subgroup of \(GL_a(\mathbf Q)\) has order dividing
\[
 \prod_p p^{\sum_{t\geq0}\lfloor a/(p^t(p-1))\rfloor}.
\]
For every prime \(p\), its displayed exponent is at most
\(\sum_{t\geq0}\lfloor2a/p^{t+1}\rfloor=v_p((2a)!)\).
Hence its order is at most \((2a)!\leq(2a)^{2a}\).
We use this classical theorem in its rational representation
form.\cite[Theorem 1]{serre-finite-groups}
It follows that the number of the prescribed isometries is
at most
\[
                 M_s=(4sf)^{4sf},\qquad M_0=1.
 \tag{EI11}
\]
The group is finite because it preserves a positive definite
form and a lattice.  No bound with a hidden degree-dependent
constant has been used.

\begin{theorem}[The two original pointwise isometry counts]
\label{thm:embedded-isometry-fixed-space-count}
Put \(X=xx'\), \(Y=XD^2\), \(m=3(2j+1)\), and
\[
 \mathcal T=\mathcal T_j(x,x';\Lambda,g).
\]
For a subspace \(U\) of codimension \(s\), put
\(n=\Delta(\Lambda\cap U)\).  The number of triples
with \(\ker(\eta-1)=U\) is zero if \(n>DX\), and otherwise
is at most
\[
 \left((2/g)^{2sf}\frac{DX}{n}\right)^{k/2}
           X^{2s+2c+1}Y^m M_s.
 \tag{EI12}
\]
The number with \(\ker(\eta-1)=0\) is also bounded by
\[
 X^{k/2}Y^{c+m}
       \left((2/g)^{kf}(4kf)^{4f}DX\right)^k .
 \tag{EI13}
\]
\end{theorem}

\begin{proof}
For (EI12), write \(Q=(2/g)^{2sf}DX\).  Lemma
\ref{lem:embedded-isometry-fixed-discriminant} gives a
common intersection \(A=F(d)\cap U=F(d')\cap U\).
If \(z=[\Lambda\cap U:A]\), then
\[
             nz^2=\Delta(A)\leq Q,\qquad
                         z\leq\min(x,x').
 \tag{EI14}
\]
In particular the set is empty if \(Q<n\), and also
if \(n>DX\), because \(g\geq2\).

Assume \(Q\geq n\) and put \(Z=\sqrt{Q/n}\).
For each possible \(z\), the number of \(A\)'s is at
most \(z^{\dim U+\ell}\leq Z^kz^\ell\).
There are at most \(\min(x,x')\leq\sqrt X\) possible
positive integers \(z\), each at most \(\sqrt X\).
Consequently there are at most
\[
                             Z^kX^{c/2}
 \tag{EI15}
\]
choices for \(A\).

For fixed \(A\), (EI10) bounds the choices for
\(F(d),F(d')\) by \(X^{s+\ell}\).  Equation (EI7)
bounds their embedded data by
\((xD)^m(x'D)^m=Y^m\).  Finally (EI11) bounds their
compatible isometries by \(M_s\).
The resulting product is
\[
 Z^kX^{s+\ell+c/2}Y^mM_s
                     \leq Z^kX^{2s+2c+1}Y^mM_s,
\]
since \(X\geq1\) and \(c=1+\ell\).  This is (EI12).
All zero-dimensional cases use the stated index and
discriminant conventions and \(M_0=1\).

For (EI13), the conclusion is immediate if there is
no triple without a fixed vector.  Otherwise (EI3)
gives
\[
                          (2/g)^{2kf}DX\geq1.
 \tag{EI16}
\]
Equations (EI8) and (EI11) bound all pairs of embedded
data with a compatible isometry by
\[
                              X^{k+c}Y^mM_k.
 \tag{EI17}
\]
Write \(a=(2/g)^{kf}\).  The ratio of the right side
of (EI13) to (EI17) equals
\[
                 a^kD^{k+2c}X^{k/2}.
\]
By (EI16), \(a^kD^{k/2}X^{k/2}\geq1\); the remaining
factor \(D^{k/2+2c}\) is at least one.  This proves
(EI13).
\end{proof}

\begin{theorem}[The original genus mean of embedded isometries]
\label{thm:embedded-isometry-genus-mean}
Retain the hypotheses and notation of
Theorem~\ref{thm:normalized-primitive-sublattice-genus-mean},
in particular \(k\geq3\), \(\rho=2v\geq2\),
\(D=T^{2f}\), and \(M_b=(20f\rho^2)^{2f(6+c)}\).
For integers \(g\geq2\), \(j\geq0\), \(x,x'\geq1\),
put \(X=xx'\) and \(Y=XD^2\).  Then
\[
 \begin{split}
 &\sum_{[\Lambda]\in\operatorname{gen}(L_0)}
       \frac{|\mathcal T_j(x,x';\Lambda,g)|}
            {|\operatorname{Aut}(\Lambda)|}\\
 &\quad\leq
 \operatorname{Mass}(L_0)\,X^{k/2}Y^{6j+10+4c}
 \left[
 1+M_b\rho^f\sum_{s=1}^{k-1}
       \left((32/g)^{kf}(4kf)^{4f}Y^7\right)^s
       +\left((2/g)^{kf}(4kf)^{4f}DX\right)^k
 \right].
 \end{split}
 \tag{EI18}
\]
This is the source's complete genus-triple estimate, with
its original field, lattices, weights and exponents.
\end{theorem}

\begin{proof}
The dual quotient of each lattice in the genus is
isomorphic to \(\mathcal O_E/T\mathcal O_E\): the
local quotients have this form, and finite primary
components combine by the Chinese remainder theorem.
Thus every lattice in the sum satisfies the hypotheses
of the pointwise estimates with discriminant \(D\).

Separate triples according to the dimension of
\(\ker(\eta-1)\).  For \(\eta=1\), apply (EI12) with
\(s=0\), \(U=V\), and \(n=D\).  The result is at most
\[
 X^{k/2}X^{2c+1}Y^{6j+3}
                   \leq X^{k/2}Y^{6j+10+4c}.
\]
After summing the reciprocal automorphism weights,
this is the first term in the bracket of (EI18).
For fixed space zero, (EI13) supplies its last term;
its factor \(Y^{c+6j+3}\) is at most
\(Y^{6j+10+4c}\).

For \(1\leq s<k\), group the fixed subspaces \(U\) by
\(n=\Delta(\Lambda\cap U)\).  Equation (EI12) shows
that only \(1\leq n\leq DX\) can occur.  Multiply
(EI12) by the normalized primitive-sublattice mean
(PG13) for rank \(k-s\).  Its two \(n/D\) factors
cancel:
\[
 \left((2/g)^{2sf}\frac{DX}{n}\right)^{k/2}
                (n/D)^{k/2}16^{skf}
                         =(32/g)^{skf}X^{k/2}.
\]
Since \(nD\leq Y\), \(X\leq Y\), and there are at
most \(DX\leq Y\) values of \(n\), the remaining
powers of \(Y\) total
\[
 3(2j+1)+(2s+2c+1)+(5s+5+2c)+1
                          =6j+10+4c+7s.
\]
Also \(M_s=(4sf)^{4sf}\leq(4kf)^{4sf}\).
The contribution of codimension \(s\) is therefore
at most
\[
 \operatorname{Mass}(L_0)X^{k/2}Y^{6j+10+4c}
 M_b\rho^f
       \left((32/g)^{kf}(4kf)^{4f}Y^7\right)^s.
\]
Summing \(s=1,\ldots,k-1\) proves (EI18).
All sums are finite by the pointwise data counts and
the finiteness of the definite genus.
\end{proof}

The remaining fixed-label return obligation is the original
weighted correspondence from pairs of flag terms to these
embedded triples, followed by its prescribed parameter bounds.
Theorem~\ref{thm:embedded-isometry-genus-mean} supplies the
arithmetic estimate on the triple side.  It does not replace
that correspondence or identify its stabilizer factors.
